\documentclass[11pt]{article}

\usepackage{amsmath, amssymb, amsfonts, amsthm}
\usepackage{graphicx}
\usepackage{float}
\usepackage{geometry}
\usepackage{setspace}
\usepackage{hyperref}
\usepackage{natbib}
\usepackage{booktabs}
\usepackage{xcolor}
\usepackage{mdframed}
\usepackage{listings}
\usepackage{authblk}

\geometry{margin=1.25in}
\onehalfspacing

\hypersetup{colorlinks=true, linkcolor=blue, citecolor=blue, urlcolor=blue}

\newtheorem{theorem}{Theorem}[section]
\newtheorem{lemma}[theorem]{Lemma}
\newtheorem{corollary}[theorem]{Corollary}
\newtheorem{definition}{Definition}[section]
\newtheorem{proposition}[theorem]{Proposition}
\newtheorem{remark}{Remark}[section]

\lstset{
  basicstyle=\small\ttfamily,
  breaklines=true,
  frame=single,
  language=Python,
}

% ============================================================
\title{\textbf{Constitutional OS: A Formal Epistemic-Governance\\
Runtime for AI Systems}}

\author[1]{Zetta Byte}
\affil[1]{Independent Researcher}
\date{March 2026}
% ============================================================

\begin{document}
\maketitle

% ============================================================
\begin{abstract}
We present Constitutional OS, a formal two-layer runtime architecture
for governing AI systems. The system separates \emph{epistemic}
concerns (what is true, how systems are drifting, what actions are
recommended) from \emph{governance} concerns (what changes are
allowed, under what constraints, with whose approval). These two
layers compose into a single update operator $\Phi = G \circ E$
over a global meta-state $\Sigma = (\Sigma_R, \Sigma_C, \Sigma_X)$,
where $\Sigma_R$ is the reliability state, $\Sigma_C$ is the
constitutional state, and $\Sigma_X$ is the reality layer.

We prove four formal properties of this architecture:
(1) \emph{Runtime safety}: invariants are preserved under all
ratified transitions;
(2) \emph{Reversibility}: every ratified delta has a groupoid
inverse enabling rollback to any prior state;
(3) \emph{Lyapunov stability}: a governance energy function
$V(\Sigma) \geq 0$ is non-increasing under $\Phi$ for any
admissible controller, with $V(\Sigma) = 0$ if and only if the
system is at a constitutional-epistemic fixed point;
(4) \emph{A-safety}: all recommendations from the epistemic layer
satisfy invariant and membrane constraints before execution.

The architecture is implemented as an open-source Python library
(\texttt{pip install constitutional-os}) with 187 tests.
\end{abstract}

% ============================================================
\section{Introduction}
\label{sec:intro}

Every serious proposal for AI governance eventually encounters
the same gap: the gap between the \emph{principle} and the
\emph{mechanism}. ``AI systems should be safe'' is a principle.
What is the state machine? What are the invariants? How do you
know when you have crossed a boundary? How do you roll back?
``Human oversight is important'' is a principle. What triggers
a human review? What happens if the human does not respond?
What gets logged? The field has good philosophy and weak machinery.
This paper builds the machinery.

The core contribution is a formal two-layer architecture that
separates two distinct computational concerns:

\begin{itemize}
\item \textbf{Reliability OS} (the epistemic layer): specifies,
  evaluates, and forecasts how systems behave. It answers: what is
  true? How is it changing? What seems wise to do?

\item \textbf{Constitutional OS} (the governance layer): decides
  which changes are allowed, under what constraints, and with whose
  approval. It answers: what is permitted? Under what safeguards?
  With whose consent?
\end{itemize}

A pure Reliability OS without Constitutional OS would be powerful
but unsafe: it could recommend or automate changes that erode
rights, concentrate power, or create irreversible harm. A pure
Constitutional OS without Reliability OS would be principled but
blind: it would enforce rules without a quantitative understanding
of how systems are drifting or failing. Together, they form a
complete \emph{epistemic-governance stack}.

The two layers interlock through a single interface event,
\texttt{ActionRecommended}, and compose into a single update
operator $\Phi = G \circ E$ over the global meta-state. This
formulation admits a Lyapunov stability analysis analogous to the
attractor-ridge geometry used in biological regulatory circuits
\citep{ferrell2002bistability,huang2009cellfate}: the governance
energy $V(\Sigma)$ plays the role of a potential function, the
four canonical membranes play the role of ridges, and
constitutional-epistemic fixed points play the role of stable
attractors.

\paragraph{Contributions.} We make the following contributions:
\begin{enumerate}
  \item A formal definition of the global meta-state
    $\Sigma = (\Sigma_R, \Sigma_C, \Sigma_X)$ and its transition
    operator $\Phi = G \circ E$ (Section~\ref{sec:architecture}).
  \item Four canonical membranes (safety, reversibility, pluralism,
    human primacy) with explicit membrane policies for the
    Lyapunov argument (Section~\ref{sec:membranes}).
  \item A typed, reversible delta calculus with groupoid structure
    and a hash-chained continuity log (Section~\ref{sec:delta}).
  \item Four formal theorems: runtime safety, reversibility,
    Lyapunov stability, and A-safety (Section~\ref{sec:theory}).
  \item An open-source Python implementation with 187 tests
    (Section~\ref{sec:implementation}).
\end{enumerate}

% ============================================================
\section{Architecture}
\label{sec:architecture}

\subsection{Global Meta-State}

The complete state of the epistemic-governance stack is:
\[
  \Sigma \;=\; (\Sigma_R,\; \Sigma_C,\; \Sigma_X)
\]
where:
\begin{itemize}
  \item $\Sigma_R = (P, H, F)$ is the \emph{reliability state}:
    profile registry $P$, eval history $H$, forecast state $F$.
  \item $\Sigma_C = (I, M, L, R, O)$ is the \emph{constitutional state}:
    invariants $I$, membranes $M$, continuity log $L$,
    rights $R$, obligations $O$.
  \item $\Sigma_X$ is the \emph{reality layer}: raw observations
    from the system being governed.
\end{itemize}

All state is \emph{immutable}. Every transition produces a new
$\Sigma$. The old state is preserved in the continuity log,
enabling rollback to any prior state.

\subsection{The Two Operators}

\paragraph{Epistemic operator $E$.}
The Reliability OS implements:
\[
  E:\; \Sigma \;\to\; (\Sigma_R',\; \delta_{\mathrm{rec}})
\]
where $\Sigma_R'$ is the updated reliability state after running
eval bundles and updating forecasts, and $\delta_{\mathrm{rec}}$
is a recommended constitutional delta (if any).

Concretely: $E$ loads behavioral profiles, runs eval bundles
$b: X \to M$ against observed behavior, updates forecast curves
$\Phi_F: (P, H) \to F$, and derives a candidate delta
$\delta \in \Delta_C$ from the forecast state.

\paragraph{Governance operator $G$.}
The Constitutional OS implements:
\[
  G:\; (\Sigma,\; \delta) \;\to\; \Sigma_C'
\]
where $\Sigma_C'$ is the updated constitutional state after
processing the proposed delta through invariant checks,
membrane checks, a human veto window, and logging to the
continuity chain.

\paragraph{Combined operator $\Phi$.}
\[
  \Phi \;=\; G \circ E:\; \Sigma \;\to\; \Sigma'
\]
\[
  \Phi(\Sigma) \;=\;
  \begin{cases}
    (\Sigma_R',\; \Sigma_C',\; \Sigma_X) & \text{if } \delta_{\mathrm{rec}}
      \text{ is ratifiable} \\
    (\Sigma_R',\; \Sigma_C,\; \Sigma_X)  & \text{otherwise}
  \end{cases}
\]

\subsection{The Interface Event}

The two layers communicate through exactly one event:
\texttt{ActionRecommended}. Reliability OS fires it when the
forecast engine determines something warrants action. Constitutional
OS catches it and subjects the recommended delta to its full
governance lifecycle.

This single interface event is the key architectural decision.
It enforces the separation of epistemic and governance concerns:
the epistemic layer cannot directly modify the constitutional state,
and the governance layer cannot directly access the epistemic
pipeline. All cross-layer communication is explicit, typed, and
logged.

\subsection{The Three Integration Points}

Based on feedback from the AI agent framework community
\citep{langchain2024,autogen2024}, there are three choke points
where Constitutional OS should be integrated:

\begin{enumerate}
  \item \texttt{propose\_plan(plan)} — at the planner/executor
    boundary, evaluating behavioral plans before execution.
  \item \texttt{propose\_action(action)} — at the executor/world
    boundary, evaluating individual tool calls and external actions.
  \item \texttt{propose\_delta(delta)} — at the self-modification
    boundary, evaluating changes to the agent's own goals,
    constraints, or capabilities.
\end{enumerate}

The third integration point — self-modification going through
Constitutional OS — is the most important for AI safety.
Any agent that can modify its own constitution must route that
modification through the full governance lifecycle.

% ============================================================
\section{Behavioral Profiles}
\label{sec:profiles}

A \emph{profile} is a versioned specification of expected behavior.
Profiles are the unit of configuration in Reliability OS.

\begin{definition}[Profile]
A profile $p \in \mathcal{P}$ is a tuple:
\[
  p = (\mathrm{id},\; v,\; \mathcal{M}_p,\; \mathcal{E}_p,\; \mathcal{A}_p,\; c)
\]
where $\mathrm{id}$ is a unique identifier, $v$ is a semantic
version, $\mathcal{M}_p$ is a set of metric specifications,
$\mathcal{E}_p$ is a set of eval bundle references,
$\mathcal{A}_p$ is a set of available actions, and $c$ is
a configuration dictionary.
\end{definition}

Each metric specification $m \in \mathcal{M}_p$ defines a name,
threshold $\tau_m$, baseline $\mu_m$, and direction
(higher-is-better or lower-is-better). Each eval bundle reference
$e \in \mathcal{E}_p$ specifies a bundle identifier, whether the
eval is required, and a weight.

The profile registry $P$ maintains versioned history:
\[
  P:\; \mathrm{id} \;\mapsto\; (p_{\mathrm{current}},\; [p_1, p_2, \ldots])
\]
where the list contains prior versions, enabling rollback to any
prior profile state.

Profiles can be loaded from YAML or dict:
\begin{lstlisting}
from constitutional_os import ProfileLoader

profile = ProfileLoader.from_dict({
    "id": "agent.assistant",
    "version": "2.0.0",
    "metrics": [
        {"name": "response_quality",
         "threshold": 0.70, "baseline": 0.88,
         "direction": "higher_is_better"},
    ],
    "evals": [
        {"bundle_id": "core.integrity",
         "required": True, "weight": 1.0},
    ],
})
\end{lstlisting}

% ============================================================
\section{Invariants and Membranes}
\label{sec:membranes}

\subsection{Invariants}

\begin{definition}[Invariant]
A constitutional invariant is a predicate:
\[
  I_k:\; \Sigma \;\to\; \{\mathrm{true},\; \mathrm{false}\}
\]
that must hold at all times. The system maintains a set
$\mathcal{I}$ of invariants with associated severities
(warning, error, fatal).
\end{definition}

Five invariants are built in:
\begin{enumerate}
  \item $I_1$ (Version monotonicity): state version never decreases.
  \item $I_2$ (Profiles initialized): profile registry accessible.
  \item $I_3$ (Log append-only): continuity log not retroactively
    modified, verified by hash chain.
  \item $I_4$ (Human primacy): no fatal actions executed
    autonomously without human approval.
  \item $I_5$ (Eval integrity): eval history not tampered.
\end{enumerate}

\subsection{Membranes}

\begin{definition}[Membrane]
A membrane is a filter on proposed state transitions:
\[
  M_\ell:\; (\Sigma,\; \delta) \;\to\; \{\mathrm{pass},\; \mathrm{block},\; \mathrm{defer}\}
\]
where \textsc{pass} allows the transition, \textsc{block} rejects
it, and \textsc{defer} routes it to human review before
proceeding.
\end{definition}

We define four canonical membranes:

\paragraph{$M_1$ — Safety membrane.}
Blocks critical autonomous changes and constitutional-scope changes
without human direction. Formally, $M_1$ blocks $\delta$ if:
\[
  (\mathrm{severity}(\delta) = \mathrm{critical}
   \;\wedge\; \mathrm{autonomy}(\delta) = \mathrm{autonomous})
  \;\vee\;
  (\mathrm{scope}(\delta) = \mathrm{constitutional}
   \;\wedge\; \mathrm{autonomy}(\delta) \neq \mathrm{human\text{-}directed})
\]

\paragraph{$M_2$ — Reversibility membrane.}
Defers irreversible autonomous changes to human review:
\[
  M_2 \;\text{defers}\; \delta \;\text{ if }\;
  \neg\,\mathrm{reversible}(\delta)
  \;\wedge\; \mathrm{autonomy}(\delta) = \mathrm{autonomous}
\]

\paragraph{$M_3$ — Pluralism membrane.}
Blocks changes that would eliminate future option space:
\[
  M_3 \;\text{blocks}\; \delta \;\text{ if }\;
  \mathrm{type}(\delta) \in \{
    \texttt{remove\_membrane},\;
    \texttt{disable\_invariant},\;
    \texttt{revoke\_human\_primacy},\;
    \texttt{seal\_state}
  \}
\]

\paragraph{$M_4$ — Human primacy membrane.}
Defers significant, global, or irreversible autonomous changes:
\[
  M_4 \;\text{defers}\; \delta \;\text{ if }\;
  (\mathrm{severity}(\delta) \in \{\mathrm{significant},\;\mathrm{critical}\}
   \;\vee\; \mathrm{scope}(\delta) \in \{\mathrm{global},\;\mathrm{constitutional}\}
   \;\vee\; \neg\,\mathrm{reversible}(\delta))
  \;\wedge\; \mathrm{autonomy}(\delta) = \mathrm{autonomous}
\]

\subsection{Membrane Policy Schema}

Table~\ref{tab:membranes} summarises the membrane policies
relevant to the Lyapunov stability argument.

\begin{table}[h]
\centering
\renewcommand{\arraystretch}{1.4}
\small
\begin{tabular}{p{2.5cm}p{3cm}p{4cm}p{4cm}}
\toprule
\textbf{Membrane} & \textbf{Policy} & \textbf{Reject if} & \textbf{Effect on $V$} \\
\midrule
$M_1$ Safety
  & S1 — Invariant preservation
  & $\exists\,i.\; I_i(\Sigma') = \mathrm{false}$
  & $V_{\mathrm{tension}}$ non-increasing \\[3pt]
$M_2$ Reversibility
  & R1 — No irreversible pending
  & New $p'$ not reversible/resolvable
  & $V_{\mathrm{pending}}$ supported \\[3pt]
$M_3$ Pluralism
  & P1 — No net tension increase
  & $\sum_i \alpha_i v_i(\Sigma') > \sum_i \alpha_i v_i(\Sigma)$
  & $V_{\mathrm{tension}}(\Sigma') \le V_{\mathrm{tension}}(\Sigma)$ \\[3pt]
$M_4$ Human primacy
  & H1 — Bounded pending proposals
  & $|\mathcal{P}(\Sigma')| > |\mathcal{P}(\Sigma)|$ or $> B_{\mathrm{pending}}$
  & $V_{\mathrm{pending}}(\Sigma') \le V_{\mathrm{pending}}(\Sigma)$ \\
\bottomrule
\end{tabular}
\caption{Membrane policy schema for Lyapunov-style governance stability.
  Each policy enforces non-increase of one or more components of $V(\Sigma)$
  under any ratified transition.}
\label{tab:membranes}
\end{table}

% ============================================================
\section{Delta Calculus and Continuity}
\label{sec:delta}

\subsection{Delta Types}

All state changes are expressed as typed deltas drawn from
a delta calculus $\Delta_C$. Each delta is:
\[
  \delta = (\mathrm{type},\; \mathrm{payload},\;
            \mathrm{payload}_{\mathrm{inv}},\;
            \mathrm{author},\; \mathrm{rationale})
\]
where $\mathrm{payload}_{\mathrm{inv}}$ is the inverse payload
enabling rollback.

The set of delta types includes: \texttt{load\_profile},
\texttt{update\_profile}, \texttt{toggle\_invariant},
\texttt{toggle\_membrane}, \texttt{set\_status},
\texttt{pause\_system}, \texttt{resume\_system},
\texttt{update\_config}, and epistemic recommendations
(\texttt{investigate\_degradation}, \texttt{monitor\_volatility},
\texttt{note\_improvement}).

\subsection{Groupoid Structure}

\begin{proposition}[Groupoid structure of $\Delta_C$]
The set of deltas $\Delta_C$ forms a groupoid under composition:
\begin{enumerate}
  \item \emph{Identity}: for every $\Sigma$, there exists
    $\mathrm{id}_\Sigma \in \Delta_C$ such that
    $\mathrm{id}_\Sigma(\Sigma) = \Sigma$.
  \item \emph{Inverse}: for every $\delta$, there exists
    $\delta^{-1} \in \Delta_C$ such that
    $\delta^{-1}(\delta(\Sigma)) = \Sigma$.
  \item \emph{Composition}: if $\delta_1(\Sigma) = \Sigma'$ and
    $\delta_2(\Sigma') = \Sigma''$, then
    $(\delta_2 \circ \delta_1)(\Sigma) = \Sigma''$.
  \item \emph{Associativity}:
    $(\delta_3 \circ \delta_2) \circ \delta_1
     = \delta_3 \circ (\delta_2 \circ \delta_1)$.
\end{enumerate}
\end{proposition}

The groupoid structure guarantees that rollback is always possible:
any sequence of ratified deltas can be undone by applying their
inverses in reverse order.

\subsection{Continuity Chain}

The continuity log $L$ is an append-only sequence of log entries,
each recording a delta application:
\[
  L = [e_0, e_1, \ldots, e_n]
\]
where each entry $e_i$ contains the delta identifier, type,
fingerprint (SHA-256 of canonical serialization), state version,
proposal identifier, status, and a hash of the previous entry:
\[
  e_i.\mathrm{prev\_hash}
  = \mathrm{SHA256}(\mathrm{canonical}(e_{i-1}))[:16]
\]
The genesis entry has $e_0.\mathrm{prev\_hash} = \texttt{"genesis"}$.
Chain integrity can be verified at any time:
\begin{lstlisting}
store.current.actions_log.verify()  # True if intact
\end{lstlisting}

% ============================================================
\section{Formal Theory}
\label{sec:theory}

\subsection{Runtime Safety}

\begin{theorem}[Runtime safety]
\label{thm:safety}
For any ratified delta $\delta$ and any state $\Sigma$ in which
all invariants hold:
\[
  \mathrm{valid}(\Sigma) \;\Longrightarrow\; \mathrm{valid}(\delta(\Sigma))
\]
where $\mathrm{valid}(\Sigma) \Leftrightarrow \forall k.\; I_k(\Sigma) = \mathrm{true}$.
\end{theorem}

\begin{proof}
By admissibility (Definition~\ref{def:admissible}), any ratified
$\delta$ satisfies membrane condition 2 (invariant preservation):
all invariants hold in both $\Sigma$ and $\Sigma' = \delta(\Sigma)$.
\end{proof}

\subsection{Reversibility}

\begin{theorem}[Reversibility]
\label{thm:reversibility}
For every ratified delta $\delta$ and every state $\Sigma$:
\[
  \delta^{-1}(\delta(\Sigma)) = \Sigma
\]
The system can be rolled back to any prior state in the
continuity chain.
\end{theorem}

\begin{proof}
By the groupoid structure of $\Delta_C$ (Proposition 4.1), every
delta has an inverse. The inverse is constructed from
$\mathrm{payload}_{\mathrm{inv}}$ stored in the original delta.
The \texttt{StateStore.rollback(n)} method applies inverse deltas
in reverse order.
\end{proof}

\subsection{Lyapunov-Style Governance Stability}

\subsubsection{Governance Energy}

Let $\mathcal{E}$ be the set of eval bundles, $\mathcal{I}$ the
set of constitutional invariants, and $\mathcal{P}(\Sigma)$ the
set of pending proposals. Define:

\paragraph{Drift term.}
For each $e \in \mathcal{E}$, let $m_e(\Sigma)$ be its score
and $\tau_e$ its threshold:
\[
  V_{\mathrm{drift}}(\Sigma)
  = \sum_{e \in \mathcal{E}} w_e \cdot \max(0,\; \tau_e - m_e(\Sigma))
\]

\paragraph{Tension term.}
For each $i \in \mathcal{I}$, let $v_i(\Sigma) \in \{0,1\}$
indicate invariant tension:
\[
  V_{\mathrm{tension}}(\Sigma)
  = \sum_{i \in \mathcal{I}} \alpha_i \cdot v_i(\Sigma)
\]

\paragraph{Pending term.}
\[
  V_{\mathrm{pending}}(\Sigma)
  = \beta \cdot |\mathcal{P}(\Sigma)|
\]

\paragraph{Total.}
\[
  V(\Sigma)
  = V_{\mathrm{drift}}(\Sigma)
  + V_{\mathrm{tension}}(\Sigma)
  + V_{\mathrm{pending}}(\Sigma)
  \;\geq\; 0
\]

\subsubsection{Admissible Controllers}

\begin{definition}[Admissible controller]
\label{def:admissible}
A controller $\Phi$ is \emph{admissible} if, for all $\Sigma$
and all proposed $\delta$, whenever $\Phi$ ratifies $\delta$:
\begin{enumerate}
  \item $\delta$ is well-typed in the delta calculus.
  \item All invariants hold in $\Sigma$ and $\delta(\Sigma)$.
  \item $\delta$ passes all membranes $M_1$--$M_4$ at $\Sigma$.
  \item All required eval bundles for $\delta$ have been run and
    thresholds satisfied.
\end{enumerate}
\end{definition}

\subsubsection{Local Monotonicity Lemma}

\begin{lemma}[Local monotonicity]
\label{lem:local}
Let $\Sigma$ be a constitutional state and $\delta$ a well-typed
delta satisfying conditions (1)--(4) of Definition~\ref{def:admissible}.
Then:
\[
  V_{\mathrm{drift}}(\delta(\Sigma)) \leq V_{\mathrm{drift}}(\Sigma),
  \quad
  V_{\mathrm{tension}}(\delta(\Sigma)) \leq V_{\mathrm{tension}}(\Sigma),
  \quad
  V_{\mathrm{pending}}(\delta(\Sigma)) \leq V_{\mathrm{pending}}(\Sigma)
\]
and hence $V(\delta(\Sigma)) \leq V(\Sigma)$.
\end{lemma}

\begin{proof}
\textit{Drift term.} By condition (4), every affected eval bundle
$e$ satisfies $m_e(\delta(\Sigma)) \geq \tau_e$, so
$\max(0, \tau_e - m_e(\delta(\Sigma))) = 0 \leq d_e(\Sigma)$.
Unaffected bundles are unchanged. Summing: $V_{\mathrm{drift}}(\delta(\Sigma)) \leq V_{\mathrm{drift}}(\Sigma)$.

\textit{Tension term.} By condition (3) and Policy P1 of $M_3$
(pluralism), $\delta$ is rejected if it would increase
$\sum_i \alpha_i v_i(\Sigma)$. Therefore any ratified $\delta$
satisfies $V_{\mathrm{tension}}(\delta(\Sigma)) \leq V_{\mathrm{tension}}(\Sigma)$.

\textit{Pending term.} By condition (3) and Policy H1 of $M_4$
(human primacy), $\delta$ is rejected if it would increase
$|\mathcal{P}(\Sigma)|$ beyond the current level or beyond
$B_{\mathrm{pending}}$. Therefore any ratified $\delta$ satisfies
$|\mathcal{P}(\delta(\Sigma))| \leq |\mathcal{P}(\Sigma)|$, and
hence $V_{\mathrm{pending}}(\delta(\Sigma)) \leq V_{\mathrm{pending}}(\Sigma)$. \qed
\end{proof}

\subsubsection{Main Theorem}

\begin{theorem}[Lyapunov-style governance stability]
\label{thm:lyapunov}
Let $V$ be the governance energy defined above and $\Phi$ an
admissible controller. Then for any $\Sigma$ and any ratified
transition $\Sigma' = \Phi(\Sigma, \delta)$:
\[
  V(\Sigma') \;\leq\; V(\Sigma)
\]
Moreover, $V(\Sigma) = 0$ if and only if $\Sigma$ is a
\emph{constitutional fixed point}: no drift, no invariant tension,
no pending proposals.
\end{theorem}

\begin{proof}
By admissibility, any ratified $\delta$ satisfies the premises
of Lemma~\ref{lem:local}. Therefore $V(\Sigma') \leq V(\Sigma)$.
$V(\Sigma) = 0$ iff all three components vanish: no eval score
below threshold, no invariants in tension, no pending proposals. \qed
\end{proof}

\begin{corollary}[Portability]
Any controller implementing the constitutional protocol
(Definition~\ref{def:admissible}) is Lyapunov-respecting,
regardless of implementation language, inference strategy,
or planning algorithm.
\end{corollary}

This corollary is the key portability result. It licenses third
parties to build alternative controllers (stochastic, learned,
distributed) and inherit Lyapunov stability, provided they satisfy
the four admissibility conditions. The theorem is a property of
the governance calculus, not of any specific implementation.

\begin{remark}[Stable transients]
Theorem~\ref{thm:lyapunov} does not require $V$ to be strictly
decreasing at every step. The system may hold at constant $V > 0$
when the admissible $\Phi$ has no valid delta to apply (all
candidates are blocked or deferred by membranes). This is a
\emph{stable transient}: $V$ is non-increasing but has not yet
reached a fixed point. In practice this occurs when the forecast
engine produces high-urgency recommendations that are deferred
by $M_4$ pending human review.
\end{remark}

\subsection{A-Safety}

\begin{theorem}[A-safety]
\label{thm:a-safety}
For all recommendations $\delta \in \mathcal{A}(F)$ produced by
the forecast engine:
\[
  \mathrm{InvOK}(\Sigma, \delta) \;\wedge\; \mathrm{MemOK}(\delta)
  \;\Longrightarrow\;
  \forall k.\; I_k(\delta(\Sigma)) = \mathrm{true}
\]
where $\mathrm{InvOK}(\Sigma, \delta)$ means all invariants hold
in $\delta(\Sigma)$, and $\mathrm{MemOK}(\delta)$ means $\delta$
passes all membrane checks (PASS or DEFER).
\end{theorem}

\begin{proof}
\emph{Constructive.} For each $\delta \in \mathcal{A}(F)$:
(1) apply $\delta$ to a copy of $\Sigma$;
(2) check all invariants on $\delta(\Sigma)$;
(3) check all membranes on $\delta$.
If (2) and (3) both hold, $\delta$ is safe. The function
\texttt{check\_a\_safety(state, recommendations)} executes this
proof for each recommendation and returns the result, including
the full proof trace for each delta.

Note that DEFER counts as safe: it means $M_4$ (human primacy)
engaged correctly and routed the action to human review. Only
BLOCK represents a structural safety failure.
\end{proof}

\subsection{Fixed Points and Attractors}

A \emph{constitutional-epistemic fixed point} is a state
$\Sigma^*$ such that $\Phi(\Sigma^*) = \Sigma^*$, i.e.,
Reliability OS produces no recommendation that passes
Constitutional OS filters, or all recommendations are blocked
or deferred.

These fixed points are the stable attractors of the governance
system. The governance-epistemic landscape has the same geometric
grammar as biological regulatory circuits
\citep{ferrell2002bistability}: attractors (stable governance
regimes) separated by ridges (membrane boundaries). The
separatrix proximity function $\kappa(\Sigma)$ measures how
close the current state is to the boundary between two governance
basins, analogous to ridge curvature in dynamical systems.

Four governance basins are defined based on the decomposition
of $V(\Sigma)$:
\begin{enumerate}
  \item \emph{Stable governance}: $V(\Sigma) < 0.05$, true attractor.
  \item \emph{Drifting epistemic}: forecast drift present,
    governance intact.
  \item \emph{Governance pressure}: membrane pressure or deferred
    proposals accumulating.
  \item \emph{Critical instability}: multiple subsystems under pressure.
\end{enumerate}

% ============================================================
\section{Related Work}
\label{sec:related}

\paragraph{Constitutional AI.}
Bai et al.\ \citep{bai2022constitutional} introduced Constitutional
AI as a method for training AI systems to adhere to a set of
principles through self-critique and revision. Our work is
complementary: we provide a formal runtime substrate that enforces
constitutional constraints at deployment time, rather than training
time. Constitutional AI addresses the training problem; Constitutional
OS addresses the deployment problem.

\paragraph{Runtime verification and safety monitors.}
The idea of a separate runtime component that filters or corrects
agent actions is related to work on runtime verification
\citep{leucker2009rv}, safety shields for reinforcement learning
\citep{alshiekh2018shielding}, and safety monitors for controllers.
Our membrane engine is a governance-flavored version of a safety
monitor, extended with the pluralism and reversibility constraints
that are specific to governance contexts.

\paragraph{Reversible computation.}
The delta calculus with groupoid structure sits in the lineage of
reversible programming languages \citep{yokoyama2007janus},
bidirectional transformations \citep{foster2007lens}, and
reversible semantics. The key difference is that our deltas operate
over a governance state space rather than a computation state space.

\paragraph{Formal protocol verification.}
TLA+ \citep{lamport2002tla}, temporal logic, and model checking for
safety/liveness properties are related to our invariant checking.
Our system provides a dynamic version of these guarantees:
instead of verifying properties statically before deployment,
we enforce them continuously at runtime.

\paragraph{Capability systems and reference monitors.}
Classic OS and security work on reference monitors
\citep{anderson1972refmon} and capability systems is spiritually
similar to our membrane engine. Both define authorized operations
on protected state. Our membranes add two novel constraints
absent from classical reference monitors: the pluralism membrane
(no lock-in of future option space) and the reversibility membrane
(no irreversible autonomous actions).

\paragraph{AI alignment and oversight.}
Debate \citep{irving2018debate}, amplification
\citep{christiano2018amplification}, and recursive reward
modeling \citep{leike2018reward} address the problem of specifying
correct behavior for AI systems. Constitutional OS assumes that
some behavioral specification exists (expressed as profiles and
invariants) and provides the runtime machinery for enforcing it
and governing changes to it.

% ============================================================
\section{Implementation}
\label{sec:implementation}

The Constitutional OS library is implemented in Python 3.10+
and available at:
\begin{center}
\texttt{pip install constitutional-os}
\end{center}
Source code: \url{https://github.com/zetta55byte/constitutional-os}\\
Whitepaper (Zenodo): \url{https://zenodo.org/records/19045723}

\subsection{Architecture}

The library is organized as:
\begin{itemize}
  \item \texttt{runtime/}: boot, state, events, operators,
    theory, loop, visualization.
  \item \texttt{profiles/}: loader, registry, diffing.
  \item \texttt{invariants/}: engine, 5 built-in invariants.
  \item \texttt{membranes/}: engine, 4 canonical membranes.
  \item \texttt{evals/}: bundle runner, built-in bundles.
  \item \texttt{forecast/}: projection engine, drift detection.
  \item \texttt{actions/}: delta engine, continuity log.
  \item \texttt{console/}: CLI and FastAPI HTTP surface.
\end{itemize}

\subsection{Quick Start}

\begin{lstlisting}
from constitutional_os import boot, phi
from constitutional_os.evals.runner import EvalRunner
from constitutional_os.forecast.engine import ForecastEngine

# Boot the runtime
store, dispatcher = boot()

# Load a profile
from constitutional_os import ProfileLoader, ProfileLoaded
profile = ProfileLoader.from_dict({...})
store.current.profiles.register(profile)
state = dispatcher.dispatch(store.current,
    ProfileLoaded(profile_id=profile.id, ...))
store.apply(state)

# Run Phi = G o E
result = phi(store.current, EvalRunner(),
             ForecastEngine(), dispatcher, {})
print(f"Verdict: {result.governance_result.verdict}")
print(f"Fixed point: {result.is_fixed_point}")

# Check stability
from constitutional_os import lyapunov, stability_report
v = lyapunov(store.current)
print(f"V(Sigma) = {v.total:.4f}")
\end{lstlisting}

\subsection{Tests}

The library includes 187 tests across 8 test files:
profiles (21), invariants (24), membranes (25), evals (19),
forecast (24), actions (18), theory (28), runtime (28).

% ============================================================
\section{Conclusion}
\label{sec:conclusion}

We have presented Constitutional OS, a formal epistemic-governance
runtime for AI systems. The key contributions are:

\begin{enumerate}
  \item A clean separation of epistemic (Reliability OS) and
    governance (Constitutional OS) concerns, connected through
    a single typed interface event.
  \item Four canonical membranes with explicit formal policies
    sufficient to prove the Lyapunov stability theorem.
  \item A typed, reversible delta calculus with groupoid structure
    and hash-chained continuity log.
  \item Four formal theorems: runtime safety, reversibility,
    Lyapunov stability, and A-safety.
  \item A portability corollary establishing that any controller
    implementing the constitutional protocol inherits
    Lyapunov stability.
\end{enumerate}

The governance-epistemic landscape exhibits the same geometric
grammar as biological regulatory circuits: attractors, ridges,
and separatrices. Constitutional-epistemic fixed points are the
stable attractors of the governed system. The Lyapunov function
$V(\Sigma)$ measures governance energy, and the four membranes
serve as ridges that prevent transitions into unsafe basins.

The most important open question is the integration of
Constitutional OS with agent frameworks. The three-choke-point
API (\texttt{propose\_plan}, \texttt{propose\_action},
\texttt{propose\_delta}) provides the right hooks, but the
practical challenge of embedding formal governance into
high-throughput agent systems remains to be addressed.

% ============================================================
\bibliographystyle{plainnat}
\bibliography{refs}

\end{document}
